\documentclass[11pt,letterpaper]{article}
\usepackage{fontspec}
\setmainfont{texgyrepagella}[Extension=.otf,UprightFont=*-regular,BoldFont=*-bold,ItalicFont=*-italic,BoldItalicFont=*-bolditalic]
\usepackage[bidi=default]{babel}
\babelprovide[import,main]{spanish}
\babelprovide[import]{hebrew}
\babelprovide[import]{ancientgreek}
\babelfont[hebrew]{rm}[Path=fonts/,BoldFont=NotoHebrewSerif-700.ttf,Scale=0.95]{NotoHebrewSerif-400.ttf}
\babelfont[ancientgreek]{rm}[Path=fonts/]{GFSPorson.otf}
\newcommand{\heb}[1]{\mbox{\foreignlanguage{hebrew}{#1}}}
\newcommand{\grc}[1]{\foreignlanguage{ancientgreek}{#1}}
\newcommand{\hebverso}[1]{\par\smallskip{\centering\begin{otherlanguage}{hebrew}\large #1\end{otherlanguage}\par}\smallskip}
\newcommand{\SENOR}{\textsc{Señor}}
\usepackage[margin=2.5cm,headheight=14pt]{geometry}
\usepackage{microtype,ragged2e,booktabs,graphicx,xcolor,csquotes,array}
\usepackage[hidelinks]{hyperref}\usepackage{xurl}\urlstyle{same}
\usepackage[most]{tcolorbox}
\usepackage{tikz}\usetikzlibrary{arrows.meta,positioning,calc,decorations.pathreplacing}
\usepackage{titlesec,fancyhdr}
\definecolor{azul}{HTML}{15437C}\definecolor{azuloscuro}{HTML}{0B1F3A}\definecolor{celeste}{HTML}{5694D6}
\definecolor{dorado}{RGB}{196,154,60}\definecolor{papiro}{RGB}{246,239,220}\definecolor{gris}{RGB}{70,78,96}
\titleformat{\section}{\Large\bfseries\color{azul}}{\thesection.}{0.6em}{}
\titleformat{\subsection}{\large\bfseries\color{azuloscuro}}{\thesubsection}{0.6em}{}
\pagestyle{fancy}\fancyhf{}
\fancyhead[L]{\foreignlanguage{spanish}{\small\itshape\color{gris}El celo de Dios}}\fancyhead[R]{\foreignlanguage{spanish}{\small\color{gris}\thepage}}
\renewcommand{\headrulewidth}{0.4pt}
\setlength{\parindent}{1.27cm}\linespread{1.15}\setlength{\emergencystretch}{3em}
\newtcolorbox{cita}[1][]{enhanced,colback=azul!6!white,colframe=azul,boxrule=0pt,leftrule=2.5pt,arc=0pt,
  left=8pt,right=8pt,top=5pt,bottom=5pt,before skip=8pt,after skip=8pt,fontupper=\itshape,#1}
\newcommand{\fuente}[1]{\par\hfill{\upshape\bfseries\small\textcolor{azul}{#1}}}
\newtcolorbox{nota}[1][]{enhanced,colback=dorado!10,colframe=dorado!85!azul,boxrule=0.6pt,arc=2pt,
  fonttitle=\bfseries\small,coltitle=azuloscuro,colbacktitle=dorado!30,fontupper=\small,title={#1}}
\newtcolorbox{filologia}[1][]{enhanced,colback=white,colframe=azul!45,boxrule=0.5pt,arc=2pt,
  fonttitle=\bfseries\small,coltitle=white,colbacktitle=azul!80,fontupper=\small,title={Filología: #1}}
\newcommand{\refapa}[1]{\par{\RaggedRight\hangindent=1.27cm\hangafter=1\noindent #1\par}\vspace{4pt}}
\newcommand{\destacar}[1]{\textbf{\textcolor{azul}{#1}}}
\newcolumntype{L}[1]{>{\RaggedRight\arraybackslash}p{#1}}


\title{\color{azul}\bfseries El celo de Dios\\[4pt]\large\color{gris}\mdseries ``\SENOR, cuyo nombre es Celoso, Dios celoso es'' (Éx 34:14)}
\author{}\date{}
\begin{document}
\maketitle\thispagestyle{empty}

\begin{abstract}\noindent
El celo de Dios incomoda porque en nosotros los celos suelen nacer de la inseguridad. Este capítulo lo estudia como un atributo de relación: supone un Dios que ama, un pueblo que le pertenece y un rival que lo desplaza. Parte de Éxodo 34:14, donde la misma palabra, \emph{qanná}, funciona primero como nombre (``Celoso es su nombre'') y luego como adjetivo (``Dios celoso es''), y la sitúa junto al Nombre de misericordia de Éx 34:6--7. Recorre el espectro de la raíz \emph{q-n-'} (85 apariciones en 70 versículos) y de la familia griega \emph{zēlos}, y muestra que, cuando Dios es sujeto, el celo es siempre \emph{por} su nombre y \emph{por} su pueblo, nunca envidia. Ordena el material con cuatro preguntas: qué es el celo, cuándo se provoca, cómo se manifiesta y cómo bendice. Estudia el encuentro de Pinjás en Baal-peor y el de Elías en Horeb, recoge la lectura judía (Mekhilta, Talmud, Rashi, Ónkelos, Maimónides, Pirqé de Rabí Eliézer) y dedica una sección a Isaías 9:7, donde ``el celo del \SENOR{} de los ejércitos'' es el sujeto que trae al Hijo de David. Termina con la lectura paulina: el celo de Cristo por la casa de su Padre, el ``celo de Dios'' con que Pablo cuida a la iglesia, el uso de Deuteronomio 32:21 en Romanos 10--11 y 1 Corintios 10:22, y un pueblo ``celoso de buenas obras''.
\par\smallskip\noindent\textbf{Palabras clave:} celo, \emph{qin'á}, \emph{qanná}, Éxodo 34:14, Isaías 9:7, Pinjás, \emph{zēlos}, Pablo.\end{abstract}

\begin{nota}[Nota sobre el texto]
Las citas bíblicas son de la Reina-Valera 1960. Donde esta escribe ``Jehová'', aquí se escribe \SENOR, que traduce el tetragrama \heb{יהוה}, leído en la tradición judía como \emph{Adonay}. El hebreo sigue la BHS; el griego, NA28 y la Septuaginta de Rahlfs-Hanhart. Las fuentes judías (Mekhilta, Talmud, Rashi, Targum, Maimónides) se citan desde Sefaria y se marcan como precisión; donde no hay traducción española publicada, la traducción es propia. Cuando la numeración hebrea difiere de la española se da entre corchetes.
\end{nota}

%===================================================================
\section{Un atributo de relación}

Pocas palabras sobre Dios suenan tan extrañas como ``celoso''. En nuestra experiencia, los celos son casi siempre un defecto: inseguridad, sospecha, deseo de controlar. El libro de Proverbios los describe sin adornos:

\hebverso{כִּי־קִנְאָה חֲמַת־גָּבֶר וְלֹא־יַחְמוֹל בְּיוֹם נָקָם}
\begin{cita}Porque los celos son el furor del hombre, Y no perdonará en el día de la venganza. No aceptará ningún rescate, Ni querrá perdonar, aunque multipliques los dones.\fuente{(Pr 6:34--35)}\end{cita}

\begin{filologia}[Pr 6:34]
\begin{tabular}{@{}L{3.2cm}L{1.6cm}L{10cm}@{}}
\heb{קִנְאָה} \emph{qin'á} & \heb{קנא} & ``celo, celos'': la pasión del marido traicionado. \\[2pt]
\heb{חֲמַת־גָּבֶר} \emph{jamat-gáver} & \heb{יחם} & ``el ardor (calor) del varón'': \emph{jemá} es calor y, por extensión, furor. \\[2pt]
\heb{לֹא־יַחְמוֹל} \emph{lo yajmol} & \heb{חמל} & ``no tendrá compasión'', no perdonará la vida. \\
\end{tabular}
\end{filologia}

El texto es una ilustración humana, no una definición de Dios: el esposo engañado no acepta rescate porque lo que está en juego es la exclusividad del pacto matrimonial. Sin embargo, la Escritura no teme usar esa palabra para Dios. Lo hace en el corazón del Decálogo y en la renovación del pacto después del becerro de oro: ``\SENOR, cuyo nombre es Celoso, Dios celoso es'' (Éx 34:14). ¿Cómo puede ser una perfección eterna lo que en nosotros suele ser pecado?

La respuesta empieza por notar que el celo es un atributo de \destacar{relación}. Necesita tres términos: alguien que ama y reclama, alguien amado que le pertenece y un rival que lo desplaza. Sin pueblo no hay celo; sin ídolo, no se enciende. Pero no es una reacción nueva en Dios: es su santo amor de siempre, visto desde el lado de la exclusividad. Por eso esta serie lo clasifica entre los atributos que se revelan en la relación con el ser humano.

\begin{nota}[Precisión]
Clasificación propuesta: atributo de Relación; incomunicable en su derecho (solo Dios puede exigir adoración exclusiva) y comunicable como reflejo (el ``celo por Dios'' de Pinjás, Elías y Pablo). Queda sujeta a la revisión del autor de la serie.
\end{nota}

El material de estudio de esta serie, preparado como sermón, ordena el atributo con cuatro preguntas: ¿qué es el celo de Dios?, ¿cuándo se provoca?, ¿cómo se manifiesta?, ¿de qué manera termina bendiciéndome? Este capítulo las sigue y les suma dos cuestiones: qué dice exactamente Éxodo cuando llama a Dios ``Celoso'' y por qué Isaías 9 atribuye al celo de Dios la venida del Mesías.

\begin{center}\small\renewcommand{\arraystretch}{1.18}
\begin{tabular}{@{}L{5.2cm}L{5.0cm}r@{}}
\toprule
\textbf{Término} & \textbf{Sentido base} & \textbf{Veces} \\ \midrule
\heb{קִנְאָה} \emph{qin'á} (sustantivo) & celo, ardor; celos; envidia & 44 \\
\heb{קִנֵּא} \emph{qinné}, \heb{הִקְנִיא} \emph{hiqní} (verbo) & tener celo; provocar a celos & 33 \\
\heb{קַנָּא} \emph{qanná} (adjetivo) & celoso (solo de Dios) & 6 \\
\heb{קַנּוֹא} \emph{qannó} (adjetivo) & celoso (solo de Dios) & 2 \\ \midrule
\grc{ζῆλος} \emph{zēlos} & celo, ardor, envidia & 16 \\
\grc{ζηλόω} \emph{zēlóō} & celar, desear con ardor, envidiar & 12 \\
\grc{ζηλωτής} \emph{zēlōtēs} & celoso; zelote & 8 \\
\grc{παραζηλόω} \emph{parazēlóō} & provocar a celos & 4 \\
\bottomrule
\end{tabular}\par
{\footnotesize\color{gris}\textit{Tabla 1.} El vocabulario del celo. \textit{Nota.} Conteo propio con la BHS y NA27 con números Strong. En la RVR1960, sin los títulos de sección, ``celo, celos, celoso, celar'' aparecen 91 veces en 78 versículos (64 en el AT y 27 en el NT).}
\end{center}

\begin{nota}[Precisión]
La raíz hebrea suma 85 apariciones en 70 versículos. Vine (1999) da 43 para el sustantivo y 34 para el verbo; el conteo propio da 44 y 33 porque una forma admite las dos lecturas. El total coincide. Conviene verificar las cifras en e-Sword.
\end{nota}

%===================================================================
\section{Texto ancla: Éxodo 20 y 34}

\subsection{El segundo mandamiento}
El adjetivo ``celoso'' aparece por primera vez en la Biblia como fundamento del segundo mandamiento:

\hebverso{לֹא־תִשְׁתַּחְוֶה לָהֶם וְלֹא תָעָבְדֵם \\ כִּי אָנֹכִי יְהוָה אֱלֹהֶיךָ אֵל קַנָּא}
\begin{cita}No tendrás dioses ajenos delante de mí. No te harás imagen, ni ninguna semejanza de lo que esté arriba en el cielo, ni abajo en la tierra, ni en las aguas debajo de la tierra. No te inclinarás a ellas, ni las honrarás; porque yo soy \SENOR{} tu Dios, fuerte, celoso, que visito la maldad de los padres sobre los hijos hasta la tercera y cuarta generación de los que me aborrecen.\fuente{(Éx 20:3--5)}\end{cita}

\begin{filologia}[Éx 20:3--5]
\begin{tabular}{@{}L{3.2cm}L{1.6cm}L{10cm}@{}}
\heb{עַל־פָּנָי} \emph{al-panay} & \heb{פנה} & ``delante de mí'', literalmente ``sobre mi rostro'': ningún otro dios en mi presencia. \\[2pt]
\heb{אֵל} \emph{El} & --- & ``Dios''; la RVR1960 traduce ``fuerte''. \\[2pt]
\heb{קַנָּא} \emph{qanná} & \heb{קנא} & adjetivo, ``celoso''; califica a \emph{El}. \\[2pt]
\heb{לְשֹׂנְאָי} \emph{lesoneay} & \heb{שׂנא} & ``de los que me aborrecen'': la visitación alcanza a quienes persisten en el odio de sus padres. \\
\end{tabular}
\end{filologia}

El mandamiento no prohíbe solo adorar a otros dioses; prohíbe también representar al Dios verdadero con imágenes. La razón es la misma: Dios no comparte su gloria con nada creado. El celo es el fundamento del culto.

\subsection{``Cuyo nombre es Celoso''}
Después del becerro de oro, Moisés sube otra vez al Sinaí y Dios renueva el pacto (Éx 34). Primero proclama su Nombre: ``¡\SENOR! ¡\SENOR! fuerte, misericordioso y piadoso\ldots'' (34:6--7). Pocos versos después, advierte contra los pactos con los cananeos:

\hebverso{כִּי לֹא תִשְׁתַּחֲוֶה לְאֵל אַחֵר \\ כִּי יְהוָה קַנָּא שְׁמוֹ אֵל קַנָּא הוּא}
\begin{cita}Guárdate de hacer alianza con los moradores de la tierra donde has de entrar, para que no sean tropezadero en medio de ti. Derribaréis sus altares, y quebraréis sus estatuas, y cortaréis sus imágenes de Asera. Porque no te has de inclinar a ningún otro dios, pues \SENOR, cuyo nombre es Celoso, Dios celoso es.\fuente{(Éx 34:12--14)}\end{cita}

La pregunta del autor de esta serie fue si aquí aparecen un sustantivo y un adjetivo. El análisis del texto da esto:

\begin{filologia}[Éx 34:14]
\begin{tabular}{@{}L{3.2cm}L{1.6cm}L{10cm}@{}}
\heb{לְאֵל אַחֵר} \emph{le'El ajer} & \heb{אחר} & ``a otro dios'' (\emph{El}, en singular). \\[2pt]
\heb{קַנָּא שְׁמוֹ} \emph{qanná shemó} & \heb{קנא} & ``Celoso (es) su nombre'': \emph{qanná} funciona como \emph{nombre}. Es el adjetivo usado como sustantivo (sustantivado), un nombre propio; la RVR1960 lo escribe con mayúscula. \\[2pt]
\heb{אֵל קַנָּא הוּא} \emph{El qanná hu} & \heb{קנא} & ``Dios celoso (es) él'': \emph{qanná} funciona como \emph{adjetivo} de \emph{El}, como en Éx 20:5. \\
\end{tabular}
\end{filologia}

La intuición es correcta en cuanto a la función, con un matiz: es la misma palabra, \emph{qanná}, dos veces. En la primera opera como nombre; en la segunda, como adjetivo. El sustantivo propiamente dicho, \heb{קִנְאָה} \emph{qin'á}, ``celo'', no aparece en Éx 20 ni en Éx 34; se usa por primera vez en la ley de los celos (Nm 5:14) y, de Dios, en Nm 25:11 y Dt 29:20. Las versiones antiguas confirman la doble función. La Septuaginta traduce \grc{ὁ γὰρ κύριος ὁ θεὸς ζηλωτὸν ὄνομα, θεὸς ζηλωτής ἐστιν}, ``el Señor Dios es nombre celoso, Dios celoso es''; el Targum Ónkelos conserva las dos palabras: \heb{אֲרֵי יְיָ קַנָּא שְׁמֵיהּ אֵל קַנָּא הוּא}.

\begin{nota}[Precisión]
En la Biblia hebrea el ``nombre'' (\heb{שֵׁם} \emph{shem}) es la identidad revelada de Dios (Éx 3:15; 33:19). Que ``Celoso'' sea su nombre indica que el celo no es una emoción pasajera, sino parte de lo que Dios es en el pacto. La lectura es teológica y propia de esta serie, apoyada en el uso del texto.
\end{nota}

\subsection{Un adjetivo reservado a Dios}
Los dos adjetivos de la raíz, \emph{qanná} y su variante \emph{qannó}, aparecen ocho veces. Todas se refieren a Dios, todas van con \heb{אֵל} \emph{El} y todas están en contexto de pacto e idolatría.

\begin{center}\small\renewcommand{\arraystretch}{1.15}
\begin{tabular}{@{}L{2.4cm}L{2.6cm}L{8.6cm}@{}}
\toprule
\textbf{Texto} & \textbf{Forma} & \textbf{Contexto} \\ \midrule
Éx 20:5 & \heb{אֵל קַנָּא} & segundo mandamiento \\
Éx 34:14 (dos veces) & \heb{קַנָּא שְׁמוֹ}, \heb{אֵל קַנָּא} & renovación del pacto; ``otro dios'', Asera \\
Dt 4:24 & \heb{אֵל קַנָּא} & ``fuego consumidor'' \\
Dt 5:9 & \heb{אֵל קַנָּא} & el Decálogo repetido \\
Dt 6:15 & \heb{אֵל קַנָּא} & ``no andaréis en pos de dioses ajenos'' \\
Jos 24:19 & \heb{אֵל־קַנּוֹא} & renovación del pacto en Siquem; junto a ``Dios santo'' \\
Nah 1:2 & \heb{אֵל קַנּוֹא} & juicio contra Nínive \\
\bottomrule
\end{tabular}\par
{\footnotesize\color{gris}\textit{Tabla 2.} Los adjetivos ``celoso'' en la Biblia hebrea. \textit{Nota.} BHS; Vine (1999).}
\end{center}

Josué lo dice con fuerza: ``No podréis servir a \SENOR, porque él es Dios santo, y Dios celoso'' (Jos 24:19). La santidad y el celo van juntos: el celo es la santidad que no tolera rivales en la relación.

\subsection{Misericordia y celo en el mismo capítulo}
Éxodo 34 es el capítulo de la misericordia (34:6--7, estudiado en el capítulo sobre la misericordia) y también el del celo (34:14). El Dios que perdona a Israel después del becerro es el mismo que no tolerará otro becerro. La advertencia que sigue usa el lenguaje del matrimonio:

\begin{cita}Por tanto, no harás alianza con los moradores de aquella tierra; porque fornicarán en pos de sus dioses, y ofrecerán sacrificios a sus dioses, y te invitarán, y comerás de sus sacrificios; o tomando de sus hijas para tus hijos, y fornicando sus hijas en pos de sus dioses, harán fornicar también a tus hijos en pos de los dioses de ellas.\fuente{(Éx 34:15--16)}\end{cita}

\begin{filologia}[Éx 34:15]
\begin{tabular}{@{}L{3.2cm}L{1.6cm}L{10cm}@{}}
\heb{וְזָנוּ אַחֲרֵי} \emph{vezanú ajaré} & \heb{זנה} & ``y fornicarán en pos de'': el verbo de la infidelidad conyugal, aplicado al culto. \\
\end{tabular}
\end{filologia}

El pacto del Sinaí es un matrimonio; la idolatría, un adulterio. Por eso Dios es ``celoso''. No se trata de un propietario que defiende un objeto, sino de un esposo que defiende una alianza de amor.

%===================================================================
\section{La palabra: el ardor del pacto}

\subsection{La raíz: calor}
La raíz hebrea \heb{קנא} expresa ardor. Los léxicos la relacionan con el árabe \emph{qana'a}, ``enrojecer intensamente'' (de tinte): el color del rostro encendido (Brown et al., 1906). El griego \grc{ζῆλος}, que la traduce en la Septuaginta, viene de \grc{ζέω}, ``hervir''; Strong lo define como ``propiamente calor''. El español ``celo'' viene del latín \emph{zelus}, y este del griego (Real Academia Española, 2014). En las tres lenguas, el celo es calor: un amor que arde, no un cálculo frío de propiedad.

\begin{nota}[Precisión]
Vine (1999) advierte que el significado árabe ``enrojeció intensamente'' ``no se explica etimológicamente''. La relación con el calor es una pista coherente con el uso bíblico, no una prueba.
\end{nota}

\subsection{El espectro completo}
La raíz se traduce en la RVR1960 de varias maneras, y no todo ``celo'' español viene de ella.

\begin{center}\small\renewcommand{\arraystretch}{1.15}
\begin{tabular}{@{}L{3.2cm}L{3.6cm}L{3.0cm}L{4.0cm}@{}}
\toprule
\textbf{Término} & \textbf{Sentido} & \textbf{RVR1960} & \textbf{Textos} \\ \midrule
\heb{קִנֵּא} \emph{qinné} (piel) & tener celo por; celar; envidiar & celo, celar, envidia, solícito & Nm 25:13; 1 R 19:10; Gn 30:1; Jl 2:18 \\
\heb{הִקְנִיא} \emph{hiqní} (hifil) & provocar a celos & despertar, mover a celos & Dt 32:16, 21; Sal 78:58; Ez 8:3 \\
\heb{קִנְאָה} \emph{qin'á} & celo, ardor; celos; envidia & celo(s), envidia & Nm 5:14; Is 9:7; Cnt 8:6; Ec 4:4 \\
\heb{קַנָּא} \emph{qanná} & celoso (de Dios) & celoso, Celoso & Éx 20:5; 34:14; Dt 4:24 \\
\heb{קַנּוֹא} \emph{qannó} & celoso (de Dios) & celoso & Jos 24:19; Nah 1:2 \\
\bottomrule
\end{tabular}\par
{\footnotesize\color{gris}\textit{Tabla 3.} La raíz \emph{q-n-'} en la Biblia hebrea. \textit{Nota.} Conteo propio con la BHS; Chávez (1992).}
\end{center}

La RVR1960 traduce la raíz por ``envidia'' o ``envidiar'' en veinte versículos (por ejemplo, Gn 26:14; 30:1; 37:11; Sal 37:1; 73:3; Pr 3:31; 14:30); por ``enojar'' en 1 R 14:22, y por ``solícito'' en Jl 2:18. A la inversa, el ``celo'' de los animales en Gn 30:41 y 31:10 traduce \heb{יחם} (``estar en calor'', la raíz de \emph{jemá}), y el de Jer 2:24 traduce \heb{חֹדֶשׁ}, ``su mes'', el tiempo del celo de la asna montés.

\begin{nota}[Precisión]
El arameo de Esd 7:17 (\heb{תִקְנֵא}, ``comprarás'') pertenece a otra raíz (Strong H7066) y queda fuera del conteo.
\end{nota}

En el Nuevo Testamento, la familia de \grc{ζῆλος} cubre todo el espectro, del celo santo a la envidia:

\begin{center}\small\renewcommand{\arraystretch}{1.15}
\begin{tabular}{@{}L{2.8cm}L{3.6cm}L{3.6cm}L{3.8cm}@{}}
\toprule
\textbf{Término} & \textbf{Sentido} & \textbf{RVR1960} & \textbf{Textos} \\ \midrule
\grc{ζῆλος} (16 vv.) & ardor; celo por; celos & celo, celos, envidia, solicitud, hervor & Jn 2:17; Ro 10:2; 2 Co 11:2; He 10:27 \\
\grc{ζηλόω} (11 vv.) & celar; desear con ardor; envidiar & celar, tener celo, procurar, envidia & 2 Co 11:2; 1 Co 12:31; 13:4; Gá 4:17--18 \\
\grc{ζηλωτής} (8 vv.) & celoso; ``zelote'' & celoso, Zelote, anheláis & Hch 22:3; Gá 1:14; Tit 2:14; Lc 6:15 \\
\grc{παραζηλόω} (4 vv.) & provocar a celos & provocar a celos & Ro 10:19; 11:11, 14; 1 Co 10:22 \\
\bottomrule
\end{tabular}\par
{\footnotesize\color{gris}\textit{Tabla 4.} La familia de \emph{zēlos} en el Nuevo Testamento. \textit{Nota.} NA27 con números Strong; Vine (1999).}
\end{center}

La Septuaginta traduce casi siempre \emph{qaná} y \emph{qin'á} con \grc{ζηλόω} y \grc{ζῆλος}, y \emph{El qanná} con \grc{θεὸς ζηλωτής}. Así llegó el vocabulario del Sinaí al griego de Pablo. En Stg 4:5, ``nos anhela celosamente'' traduce otra palabra, \grc{πρὸς φθόνον ἐπιποθεῖ}, de interpretación discutida.

\subsection{Celo por, celo contra}
La sintaxis del verbo distingue el celo santo de la envidia. Con la preposición \heb{לְ} (\emph{le-}, ``por, a favor de'') significa tener celo \emph{por} alguien o algo; con \heb{בְּ} (\emph{be-}, ``contra'') significa envidiar; en hifil significa provocar el celo de otro.

\begin{center}\small\renewcommand{\arraystretch}{1.18}
\begin{tabular}{@{}L{2.8cm}L{3.4cm}L{7.6cm}@{}}
\toprule
\textbf{Construcción} & \textbf{Sentido} & \textbf{Ejemplos} \\ \midrule
\emph{qinné} + \heb{לְ} & celo \emph{por} & ``tuvo celo por su Dios'' (Nm 25:13); ``celé a Sion'' (Zac 8:2); ``por mi santo nombre'' (Ez 39:25); ``por su tierra'' (Jl 2:18) \\
\emph{qinné} + \heb{בְּ} & celo \emph{contra}, envidia & ``tuvo envidia de su hermana'' (Gn 30:1); ``no tengas envidia de los malignos'' (Sal 37:1; Pr 3:31) \\
\emph{qinné} + objeto & celar a (cónyuge) & ``tuviere celos de su mujer'' (Nm 5:14, 30) \\
\emph{hiqní} (hifil) & provocar el celo & ``le despertaron a celos con los dioses ajenos'' (Dt 32:16) \\
\bottomrule
\end{tabular}\par
{\footnotesize\color{gris}\textit{Tabla 5.} Construcciones del verbo \emph{qinné}. \textit{Nota.} Elaboración propia a partir de la BHS y Vine (1999).}
\end{center}

\hebverso{וַתְּקַנֵּא רָחֵל בַּאֲחֹתָהּ \\ קִנֵּאתִי לְצִיּוֹן קִנְאָה גְדוֹלָה}
\begin{cita}Viendo Raquel que no daba hijos a Jacob, tuvo envidia de su hermana.\fuente{(Gn 30:1)}\end{cita}
\begin{cita}Así ha dicho \SENOR{} de los ejércitos: Celé a Sion con gran celo, y con gran ira la celé.\fuente{(Zac 8:2)}\end{cita}

Raquel cela \emph{contra} su hermana; Dios cela \emph{por} Sion. Vine (1999) lo resume: ``En ningún pasaje de todo el Antiguo Testamento se dice que Dios siente envidia''. Cuando Dios es el sujeto, el verbo va con \heb{לְ}: su celo es a favor de su nombre, de su tierra y de su pueblo.

\subsection{Las palabras que acompañan al celo}
El celo de Dios aparece junto al fuego (\heb{אֵשׁ}: Dt 4:24; Sal 79:5; Sof 1:18; 3:8), al calor y al furor (\heb{חֵמָה}: Zac 8:2; Ez 16:38; 36:6), a la ira (\heb{אַף}: Dt 29:20), pero también al amor (Cnt 8:6) y a las entrañas (Is 63:15). Es fuego contra el rival y ternura hacia el amado. Israel lo reclama en la oración:

\hebverso{אַיֵּה קִנְאָתְךָ וּגְבוּרֹתֶךָ הֲמוֹן מֵעֶיךָ וְרַחֲמֶיךָ}
\begin{cita}Mira desde el cielo, y contempla desde tu santa y gloriosa morada. ¿Dónde está tu celo, y tu poder, la conmoción de tus entrañas y tus piedades para conmigo? ¿Se han estrechado?\fuente{(Is 63:15)}\end{cita}

\subsection{Las lenguas hermanas}
El material de estudio incluye una comparación con otras lenguas semíticas. Conviene precisar qué palabras comparten la raíz y cuáles solo la idea.

\begin{center}\small\renewcommand{\arraystretch}{1.15}
\begin{tabular}{@{}L{1.7cm}L{2.8cm}L{2.2cm}L{7.0cm}@{}}
\toprule
\textbf{Lengua} & \textbf{Término} & \textbf{¿Misma raíz?} & \textbf{Sentido} \\ \midrule
Hebreo & \heb{קנא} \emph{q-n-'} & --- & arder de celo; defender lo propio; envidiar \\
Arameo & \heb{קַנָּא} (Targum) & sí & los targumes conservan \emph{qanná} en Éx 20:5; 34:14; Dt 4:24 \\
Árabe & \emph{qana'a} & sí (cognado) & ``enrojecer intensamente'' \\
Árabe & \emph{ghayra} & no & celo por el honor y por los propios (raíz \emph{gh-y-r}) \\
Árabe & \emph{ḥamiyya} & no & fervor, ardor protector (raíz \emph{ḥ-m-y}, ``calentar, proteger'') \\
\bottomrule
\end{tabular}\par
{\footnotesize\color{gris}\textit{Tabla 6.} El celo en las lenguas semíticas. \textit{Nota.} Brown et al. (1906); Vine (1999).}
\end{center}

\begin{nota}[Precisión]
\emph{Ghayra} y \emph{ḥamiyya} no proceden de la raíz \emph{q-n-'}: sirven como paralelo cultural, no como etimología. El cognado árabe directo es \emph{qana'a}. La idea central del material sí se sostiene: en el mundo semítico el celo pertenece al honor y al pacto. El esposo defiende la alianza, el rey su gloria y Dios su pueblo y su Nombre.
\end{nota}

\subsection{Cantares 8:6: una llama de Yah}
La poesía hebrea pone el celo en paralelo con el amor:

\hebverso{כִּי־עַזָּה כַמָּוֶת אַהֲבָה קָשָׁה כִשְׁאוֹל קִנְאָה \\ רְשָׁפֶיהָ רִשְׁפֵּי אֵשׁ שַׁלְהֶבֶתְיָה}
\begin{cita}Ponme como un sello sobre tu corazón, como una marca sobre tu brazo; Porque fuerte es como la muerte el amor; Duros como el Seol los celos; Sus brasas, brasas de fuego, fuerte llama.\fuente{(Cnt 8:6)}\end{cita}

\begin{filologia}[Cnt 8:6]
\begin{tabular}{@{}L{3.2cm}L{1.6cm}L{10cm}@{}}
\heb{אַהֲבָה} // \heb{קִנְאָה} & \heb{אהב}, \heb{קנא} & paralelismo: amor // celo. El celo es el amor visto desde su exclusividad. \\[2pt]
\heb{קָשָׁה כִשְׁאוֹל} & \heb{קשׁה} & ``duro como el Seol'': no suelta lo que tiene. \\[2pt]
\heb{שַׁלְהֶבֶתְיָה} \emph{shalhévetyá} & \heb{להב} & ``llama de Yah''; la RVR1960 traduce ``fuerte llama''. \\
\end{tabular}
\end{filologia}

\begin{nota}[Precisión]
La terminación \emph{-yá} puede ser el nombre divino abreviado o un superlativo hebreo (``llama intensísima''); los intérpretes se dividen.
\end{nota}

%===================================================================
\section{El Nombre que revela el celo}

El celo es el caso más claro de un atributo que se vuelve Nombre. El texto ancla ya lo mostró: ``\SENOR{}, cuyo nombre es Celoso, Dios celoso es'' (\heb{קַנָּא שְׁמוֹ}, Éx 34:14). Como con la santidad (``cuyo nombre es el Santo'', Is 57:15), la Escritura no solo dice que Dios es celoso: dice que ese es su nombre. A esa fórmula se suman otros títulos que unen el celo con otro atributo.

\hebverso{כִּי יְהוָה אֱלֹהֶיךָ אֵשׁ אֹכְלָה הוּא אֵל קַנָּא}
\begin{cita}Porque \SENOR{} tu Dios es fuego consumidor, Dios celoso.\fuente{(Dt 4:24)}\end{cita}

\begin{filologia}[Dt 4:24]
\heb{אֵשׁ אֹכְלָה} (\emph{esh ojlá}): ``fuego que devora''. \heb{אֵל קַנָּא} (\emph{El qanná}): ``Dios celoso''. El título aparece junto a la imagen del fuego, la misma que Hebreos retoma: ``nuestro Dios es fuego consumidor'' (He 12:29).
\end{filologia}

El celo de Dios tiene también un objeto que es su propio Nombre:

\hebverso{וְרִחַמְתִּי כָּל־בֵּית יִשְׂרָאֵל וְקִנֵּאתִי לְשֵׁם קָדְשִׁי}
\begin{cita}Por tanto, así ha dicho \SENOR{} el Señor: Ahora volveré la cautividad de Jacob, y tendré misericordia de toda la casa de Israel, y me mostraré celoso por mi santo nombre.\fuente{(Ez 39:25)}\end{cita}

\begin{filologia}[Ez 39:25]
\heb{וְקִנֵּאתִי} (\emph{veqinnetí}, piel de \heb{קנא}): ``y me mostraré celoso''. \heb{לְשֵׁם קָדְשִׁי} (\emph{leshem qodshí}): ``por el nombre de mi santidad''. En un mismo verso, la misericordia (\heb{וְרִחַמְתִּי}) y el celo, ambos al servicio del Nombre santo.
\end{filologia}

El Dios cuyo nombre es Celoso cela por su Nombre. Por eso el celo no es un capricho: protege la gloria del Nombre que contiene la santidad, la misericordia y la verdad (Éx 34:6--7, 14).

\begin{center}
\small\renewcommand{\arraystretch}{1.2}
\begin{tabular}{@{}L{4.2cm}L{3.6cm}L{3.0cm}L{4.6cm}@{}}
\toprule
\textbf{Nombre o título} & \textbf{Hebreo / griego} & \textbf{Texto} & \textbf{Qué revela} \\ \midrule
``Cuyo nombre es Celoso'' & \heb{קַנָּא שְׁמוֹ} & Éx 34:14 & El atributo es el Nombre \\
``Dios celoso'' & \heb{אֵל קַנָּא} & Éx 20:5; Dt 4:24 & Celo y fuego consumidor \\
``Dios santo, Dios celoso'' & \heb{אֵל־קַנּוֹא} & Jos 24:19 & Celo y santidad \\
``Dios celoso y vengador'' & \heb{אֵל קַנּוֹא וְנֹקֵם} & Nah 1:2 & Celo e ira santa \\
``Celoso por mi santo nombre'' & \heb{וְקִנֵּאתִי לְשֵׁם קָדְשִׁי} & Ez 39:25 & El Nombre, objeto del celo \\
``El celo del \SENOR{} de los ejércitos'' & \heb{קִנְאַת יְהוָה צְבָאוֹת} & Is 9:7 [9:6] & El celo que envía al Mesías \\
``Celo de Dios'' & \grc{θεοῦ ζήλῳ} & 2~Co 11:2 & Pablo cela por la iglesia \\ \bottomrule
\end{tabular}\par\smallskip
{\footnotesize\color{gris}\textit{Tabla.} El celo en los nombres de Dios. \textit{Nota.} Elaboración propia.}
\end{center}

%===================================================================
\section{Definición}

\subsection{¿Qué es el celo de Dios?}
El material de estudio responde a la primera pregunta con cinco rasgos. El celo de Dios es: (1) un atributo esencial de su naturaleza, no una reacción insegura; (2) pasión santa y ferviente por su gloria; (3) defensa activa de lo que le pertenece; (4) exigencia de exclusividad en la relación de pacto; (5) amor protector que no tolera rivales. Y lo resume así: no es inseguridad, es santidad en acción frente a la idolatría.

A partir de esos rasgos y de los textos estudiados, esta serie propone la siguiente definición:

\begin{tcolorbox}[colback=azul!5!white,colframe=azul,boxrule=1pt,arc=3pt,title={\small Definición},colbacktitle=azul,coltitle=white]
El celo de Dios es la perfección eterna por la cual Dios, en el ardor de su santo amor de pacto, reclama para sí la devoción exclusiva de lo que le pertenece, se levanta contra todo rival que lo desplace y actúa con pasión \emph{por} su Nombre y \emph{por} su pueblo, para juicio de los ídolos y para salvación de los suyos.
\end{tcolorbox}

\subsection{Celo humano y celo divino}
La diferencia con los celos humanos no está en la intensidad, sino en el objeto, el origen y el fin.

\begin{center}\small\renewcommand{\arraystretch}{1.2}
\begin{tabular}{@{}L{2.3cm}L{5.4cm}L{6.0cm}@{}}
\toprule
 & \textbf{Envidia, celos pecaminosos} & \textbf{Celo de Dios} \\ \midrule
Objeto & desea lo que es de otro (Gn 30:1; Sal 73:3) & defiende lo que es suyo (Éx 19:5; 1 Co 6:19--20) \\
Origen & inseguridad, carencia, comparación (Stg 3:14--16) & plenitud y santidad: ``yo \SENOR; este es mi nombre'' (Is 42:8) \\
Rival & alguien mayor que yo & ``lo que no es Dios'' (Dt 32:21) \\
Fin & destruir al otro & destruir al ídolo y restaurar al amado (Ez 16:42; Os 2:19--20) \\
Dominio & ``los celos son el furor del hombre'' (Pr 6:34): lo dominan & ``Yo domino el celo; el celo no me domina'' (Mekhilta) \\
\bottomrule
\end{tabular}\par
{\footnotesize\color{gris}\textit{Tabla 7.} Celo humano y celo divino. \textit{Nota.} Elaboración propia.}
\end{center}

\hebverso{אֲנִי יְהוָה הוּא שְׁמִי וּכְבוֹדִי לְאַחֵר לֹא־אֶתֵּן}
\begin{cita}Yo \SENOR; este es mi nombre; y a otro no daré mi gloria, ni mi alabanza a esculturas.\fuente{(Is 42:8)}\end{cita}

Dios no cela porque le falte algo, sino porque es la plenitud de todo bien. Dar su gloria a una escultura sería mentir sobre lo que él es. Su celo protege la verdad.

\subsection{Lo que le pertenece}
El celo presupone pertenencia. Dios cela lo que creó, eligió y redimió:

\hebverso{וִהְיִיתֶם לִי סְגֻלָּה מִכָּל־הָעַמִּים כִּי־לִי כָּל־הָאָרֶץ}
\begin{cita}Ahora, pues, si diereis oído a mi voz, y guardareis mi pacto, vosotros seréis mi especial tesoro sobre todos los pueblos; porque mía es toda la tierra.\fuente{(Éx 19:5)}\end{cita}
\begin{cita}¿O ignoráis que vuestro cuerpo es templo del Espíritu Santo, el cual está en vosotros, el cual tenéis de Dios, y que no sois vuestros? Porque habéis sido comprados por precio; glorificad, pues, a Dios en vuestro cuerpo y en vuestro espíritu, los cuales son de Dios.\fuente{(1 Co 6:19--20)}\end{cita}

\begin{filologia}[Éx 19:5]
\begin{tabular}{@{}L{3.2cm}L{1.6cm}L{10cm}@{}}
\heb{סְגֻלָּה} \emph{segul.lá} & \heb{סגל} & ``tesoro personal'', la propiedad privada de un rey. La LXX traduce \grc{λαὸς περιούσιος}, ``pueblo de su propiedad'', la expresión que reaparece en Tit 2:14. \\
\end{tabular}
\end{filologia}

\begin{nota}[Precisión]
\heb{קָנָה} \emph{qaná} (con \emph{he}), ``comprar, adquirir'', y \heb{קָנָא} \emph{qaná'} (con \emph{alef}), ``celar'', son raíces distintas (Chávez, 1992; Strong H7065 y H7069). Suenan parecido, pero unirlas sería una falacia etimológica (Carson, 1996).
\end{nota}

%===================================================================
\section{Teología: cuándo se provoca y cómo se manifiesta}

\subsection{``Lo que no es Dios''}
La segunda pregunta del material es cuándo se provoca el celo de Dios. Su respuesta: cuando aparece un ``amante'' en el corazón; cuando algo desplaza a Dios del primer lugar; cuando se cambia la verdad por la mentira; cuando se ama más a la criatura que al Creador; cuando se divide la devoción. El cántico de Moisés lo dice con el hifil de la raíz, ``provocar a celos'':

\hebverso{הֵם קִנְאוּנִי בְלֹא־אֵל כִּעֲסוּנִי בְּהַבְלֵיהֶם \\ וַאֲנִי אַקְנִיאֵם בְּלֹא־עָם בְּגוֹי נָבָל אַכְעִיסֵם}
\begin{cita}Ellos me movieron a celos con lo que no es Dios; Me provocaron a ira con sus ídolos; Yo también los moveré a celos con un pueblo que no es pueblo, Los provocaré a ira con una nación insensata.\fuente{(Dt 32:21)}\end{cita}

\begin{filologia}[Dt 32:21]
\begin{tabular}{@{}L{3.2cm}L{1.6cm}L{10cm}@{}}
\heb{קִנְאוּנִי} \emph{qin'uni} & \heb{קנא} & ``me provocaron a celos'' (piel con sentido causativo). \\[2pt]
\heb{בְלֹא־אֵל} \emph{belo-El} & \heb{אל} & ``con un no-dios'': el rival no es grande; es nada. \\[2pt]
\heb{בְּהַבְלֵיהֶם} \emph{behavleihem} & \heb{הבל} & ``con sus vanidades'': \emph{hével}, el aliento que se esfuma (Ec 1:2). \\[2pt]
\heb{אַקְנִיאֵם} \emph{aqni'em} & \heb{קנא} & ``los provocaré a celos'' (hifil). \\[2pt]
\heb{בְּלֹא־עָם} \emph{belo-am} & \heb{עם} & ``con un no-pueblo'': respuesta simétrica de Dios. \\
\end{tabular}
\end{filologia}

Los textos del material desarrollan la misma idea. Pablo describe la idolatría como un cambio: ``cambiaron la verdad de Dios por la mentira, honrando y dando culto a las criaturas antes que al Creador'' (Ro 1:25). Jeremías la presenta como adulterio y divorcio: ``yo la había despedido y dado carta de repudio'' (Jer 3:8). Jesús la reconoce en la devoción dividida: ``Ninguno puede servir a dos señores'' (Mt 6:24), y en los afectos que compiten con él: ``El que ama a padre o madre más que a mí, no es digno de mí'' (Mt 10:37). Moisés une el olvido del pacto con el fuego del celo: ``Guardaos, no os olvidéis del pacto\ldots{} Porque \SENOR{} tu Dios es fuego consumidor, Dios celoso'' (Dt 4:23--24).

\subsection{La imagen del celo}
Ezequiel ve en visión un ídolo instalado en el templo mismo:

\hebverso{מוֹשַׁב סֵמֶל הַקִּנְאָה הַמַּקְנֶה}
\begin{cita}Y aquella figura extendió la mano, y me tomó por las guedejas de mi cabeza; y el Espíritu me alzó entre el cielo y la tierra, y me llevó en visiones de Dios a Jerusalén, a la entrada de la puerta de adentro que mira hacia el norte, donde estaba la habitación de la imagen del celo, la que provoca a celos.\fuente{(Ez 8:3)}\end{cita}

\begin{filologia}[Ez 8:3]
\begin{tabular}{@{}L{3.2cm}L{1.6cm}L{10cm}@{}}
\heb{סֵמֶל} \emph{sémel} & \heb{סמל} & ``imagen, estatua'' (cf. Dt 4:16). \\[2pt]
\heb{הַקִּנְאָה הַמַּקְנֶה} \emph{haqqin'á hammaqné} & \heb{קנא} & ``el celo que provoca celo'': juego de palabras con la misma raíz; \emph{hammaqné} es hifil escrito con \emph{he} final (Chávez, 1992). \\
\end{tabular}
\end{filologia}

Lo que despierta el celo de Dios no es una amenaza externa, sino un rival dentro de su propia casa. Jesús echará del templo a los mercaderes con el mismo celo (Jn 2:17).

\subsection{Un amante en el corazón}
Los profetas leen la historia de Israel como un matrimonio: el desposorio del Sinaí (Éx 19:5), el adulterio de la idolatría (Éx 34:15; Jer 3:8), el celo del Esposo (Ez 16:38) y un nuevo desposorio (Os 2:19--20). La meta del celo no es el divorcio, sino la restauración:

\begin{cita}Y saciaré mi ira sobre ti, y se apartará de ti mi celo, y descansaré y no me enojaré más.\fuente{(Ez 16:42)}\end{cita}
\begin{cita}Y te desposaré conmigo para siempre; te desposaré conmigo en justicia, juicio, benignidad y misericordia. Y te desposaré conmigo en fidelidad, y conocerás a \SENOR.\fuente{(Os 2:19--20)}\end{cita}

\subsection{¿Cómo se manifiesta?}
La tercera pregunta del material recibe esta respuesta: el celo de Dios se manifiesta en disciplina, juicio correctivo, purificación, defensa pública de su gloria y guerra contra los ídolos. No es exactamente ira general, sino ira dirigida contra la rivalidad idolátrica. Dos reyes lo ilustran:

\begin{cita}habló el rey y dijo: ¿No es ésta la gran Babilonia que yo edifiqué para casa real con la fuerza de mi poder, y para gloria de mi majestad?\fuente{(Dn 4:30)}\end{cita}
\begin{cita}Y un día señalado, Herodes, vestido de ropas reales, se sentó en el tribunal y les arengó. Y el pueblo aclamaba gritando: ¡Voz de Dios, y no de hombre! Al momento un ángel del Señor le hirió, por cuanto no dio la gloria a Dios; y expiró comido de gusanos.\fuente{(Hch 12:21--23)}\end{cita}

Nabucodonosor se atribuye la gloria y es humillado ``hasta que reconozcas que el Altísimo tiene el dominio en el reino de los hombres'' (Dn 4:32); al final alaba al Rey del cielo (4:37). Su caso es disciplina que restaura. Herodes acepta ser llamado dios y muere ``por cuanto no dio la gloria a Dios''. En ambos casos Dios defiende públicamente su gloria.

\subsection{Fuego consumidor}
El celo y el fuego van juntos desde Deuteronomio:

\hebverso{כִּי יְהוָה אֱלֹהֶיךָ אֵשׁ אֹכְלָה הוּא אֵל קַנָּא}
\begin{cita}Porque \SENOR{} tu Dios es fuego consumidor, Dios celoso.\fuente{(Dt 4:24)}\end{cita}
\begin{cita}Así que, recibiendo nosotros un reino inconmovible, tengamos gratitud, y mediante ella sirvamos a Dios agradándole con temor y reverencia; porque nuestro Dios es fuego consumidor.\fuente{(He 12:28--29)}\end{cita}

Pink no dedica un capítulo al celo, pero toca sus dos lados. Sobre Hebreos 12:28--29, en su capítulo sobre la ira, escribe que no podemos servir a Dios ``aceptablemente'' sin reverencia a su majestad, y que esos sentimientos ``se desarrollan mejor al recordar frecuentemente que `nuestro Dios es fuego consumidor'{}'' (Pink, 2020, cap. 18). Y en su capítulo sobre el amor advierte contra un amor divino reducido a ``una especie de amable debilidad'', ``modelado según la emoción humana'' (Pink, 2020, cap. 16). Lo mismo vale para el celo: hay que pensarlo desde la Escritura y no desde nuestros celos.

\begin{nota}[Precisión]
Pink se usa aquí como referencia puntual; el orden y la estructura del estudio son los de esta serie.
\end{nota}

\subsection{El celo que defiende a su pueblo}
El título del material de estudio, ``Celo por su pueblo'', señala el otro rostro del atributo. El mismo fuego que consume ídolos protege a Sion:

\begin{center}\small\renewcommand{\arraystretch}{1.18}
\begin{tabular}{@{}L{1.8cm}L{7.4cm}L{4.6cm}@{}}
\toprule
\textbf{Texto} & \textbf{RVR1960} & \textbf{Hebreo} \\ \midrule
Zac 1:14 & ``Celé con gran celo a Jerusalén y a Sion'' & \heb{קִנֵּאתִי לִירוּשָׁלִַם} \\
Zac 8:2 & ``Celé a Sion con gran celo, y con gran ira la celé'' & \heb{קִנְאָה גְדוֹלָה וְחֵמָה גְדוֹלָה} \\
Jl 2:18 & ``Y \SENOR, solícito por su tierra, perdonará a su pueblo'' & \heb{וַיְקַנֵּא יְהוָה לְאַרְצוֹ} \\
Ez 39:25 & ``tendré misericordia de toda la casa de Israel, y me mostraré celoso por mi santo nombre'' & \heb{וְקִנֵּאתִי לְשֵׁם קָדְשִׁי} \\
Is 59:17 & ``se cubrió de celo como de manto'' & \heb{וַיַּעַט כַּמְעִיל קִנְאָה} \\
\bottomrule
\end{tabular}\par
{\footnotesize\color{gris}\textit{Tabla 8.} El celo de Dios por su pueblo. \textit{Nota.} RVR1960 y BHS.}
\end{center}

En Ez 39:25 la misericordia y el celo están en la misma frase. En Jl 2:18 el celo de Dios por su tierra se traduce en perdón y en pan: ``He aquí yo os envío pan, mosto y aceite'' (2:19).

\begin{nota}[Precisión]
En Jl 2:18 la RVR1960 traduce ``solícito''; el hebreo dice ``tuvo celo por su tierra'' (\emph{vayqanné}). La cita ``\SENOR{} se llenará de celo'' del material de estudio no corresponde a la RVR1960.
\end{nota}

%===================================================================
\section{Un encuentro en la historia de Israel: Pinjás y Elías}

\subsection{Baal-peor}
Al final del desierto, acampado en Sitim, frente a Jericó, Israel repite el pecado del Sinaí. El pueblo ``empezó a fornicar con las hijas de Moab'' y se unió al culto de su dios:

\hebverso{וַיִּצָּמֶד יִשְׂרָאֵל לְבַעַל פְּעוֹר וַיִּחַר־אַף יְהוָה בְּיִשְׂרָאֵל}
\begin{cita}Así acudió el pueblo a Baal-peor; y el furor de \SENOR{} se encendió contra Israel.\fuente{(Nm 25:3)}\end{cita}

El verbo \heb{וַיִּצָּמֶד} (\emph{vayyitsámed}) significa ``se unció'', ``se ató'' como a un yugo. Una plaga se extiende por el campamento y mueren veinticuatro mil (Nm 25:9). Entonces Pinjás, nieto de Aarón, ``se levantó de en medio de la congregación, y tomó una lanza en su mano'' y mató al israelita y a la madianita que había traído a su tienda (25:7--8). La plaga se detuvo.

\subsection{Celoso con el celo de Dios}
Dios interpreta el acto de Pinjás con el vocabulario del celo:

\hebverso{פִּינְחָס בֶּן־אֶלְעָזָר בֶּן־אַהֲרֹן הַכֹּהֵן הֵשִׁיב אֶת־חֲמָתִי מֵעַל בְּנֵי־יִשְׂרָאֵל \\ בְּקַנְאוֹ אֶת־קִנְאָתִי בְּתוֹכָם וְלֹא־כִלִּיתִי אֶת־בְּנֵי־יִשְׂרָאֵל בְּקִנְאָתִי}
\begin{cita}Finees hijo de Eleazar, hijo del sacerdote Aarón, ha hecho apartar mi furor de los hijos de Israel, llevado de celo entre ellos; por lo cual yo no he consumido en mi celo a los hijos de Israel. Por tanto diles: He aquí yo establezco mi pacto de paz con él; y tendrá él, y su descendencia después de él, el pacto del sacerdocio perpetuo, por cuanto tuvo celo por su Dios e hizo expiación por los hijos de Israel.\fuente{(Nm 25:11--13)}\end{cita}

\begin{filologia}[Nm 25:11--13]
\begin{tabular}{@{}L{3.2cm}L{1.6cm}L{10cm}@{}}
\heb{הֵשִׁיב אֶת־חֲמָתִי} \emph{heshiv et-jamatí} & \heb{שׁוב} & ``hizo volver mi furor'': lo apartó de Israel. \\[2pt]
\heb{בְּקַנְאוֹ אֶת־קִנְאָתִי} \emph{beqan'ó et-qin'atí} & \heb{קנא} & ``al celar \emph{mi} celo'': Pinjás no tuvo un celo propio, compartió el de Dios. La RVR1960 traduce ``llevado de celo''. \\[2pt]
\heb{בְּרִיתִי שָׁלוֹם} \emph{berití shalom} & \heb{שׁלם} & ``mi pacto de paz'': el celo termina en paz. \\[2pt]
\heb{קִנֵּא לֵאלֹהָיו} \emph{qinné lelohav} & \heb{קנא} & ``tuvo celo \emph{por} su Dios'', con la preposición \heb{לְ}. \\[2pt]
\heb{וַיְכַפֵּר} \emph{vaykhappér} & \heb{כפר} & ``e hizo expiación'': cubrió el pecado del pueblo. \\
\end{tabular}
\end{filologia}

El texto es notable por tres razones. Primero, el celo de Pinjás es el celo \emph{de Dios} (``mi celo'') obrando a través de un hombre. Segundo, ese celo produce expiación: la ira de Dios se aparta y el pueblo no es consumido. Tercero, termina en un pacto de \emph{paz}. El Salmo 106 lo recuerda con una fórmula sorprendente:

\hebverso{וַיַּעֲמֹד פִּינְחָס וַיְפַלֵּל וַתֵּעָצַר הַמַּגֵּפָה \\ וַתֵּחָשֶׁב לוֹ לִצְדָקָה לְדֹר וָדֹר עַד־עוֹלָם}
\begin{cita}Entonces se levantó Finees e hizo juicio, Y se detuvo la plaga; Y le fue contado por justicia De generación en generación para siempre.\fuente{(Sal 106:30--31)}\end{cita}

\begin{nota}[Precisión]
``Le fue contado por justicia'' es casi la fórmula de Gn 15:6 sobre Abraham (\heb{וַיַּחְשְׁבֶהָ לּוֹ צְדָקָה}), que Pablo cita en Ro 4:3. La tradición judía del Segundo Templo asoció a Abraham y a Pinjás como modelos (1 Mac 2:52, 54). La conexión con Pablo se retoma en la sección paulina.
\end{nota}

\subsection{Elías en Horeb}
Siglos después, Elías huye de Jezabel hasta Horeb, el monte del pacto. Allí Dios le pregunta qué hace, y él responde con la raíz del celo duplicada:

\hebverso{קַנֹּא קִנֵּאתִי לַיהוָה אֱלֹהֵי צְבָאוֹת \\ כִּי־עָזְבוּ בְרִיתְךָ בְּנֵי יִשְׂרָאֵל}
\begin{cita}Y allí se metió en una cueva, donde pasó la noche. Y vino a él palabra de \SENOR, el cual le dijo: ¿Qué haces aquí, Elías? El respondió: He sentido un vivo celo por \SENOR{} Dios de los ejércitos; porque los hijos de Israel han dejado tu pacto, han derribado tus altares, y han matado a espada a tus profetas; y sólo yo he quedado, y me buscan para quitarme la vida.\fuente{(1 R 19:9--10)}\end{cita}

\begin{filologia}[1 R 19:10--12]
\begin{tabular}{@{}L{3.2cm}L{1.6cm}L{10cm}@{}}
\heb{קַנֹּא קִנֵּאתִי} \emph{qannó qinnetí} & \heb{קנא} & infinitivo absoluto más verbo: ``celar, he celado'', ``he celado ardientemente''. \\[2pt]
\heb{לַיהוָה אֱלֹהֵי צְבָאוֹת} & \heb{צבא} & ``por el \SENOR, Dios de los ejércitos'': el mismo título de Is 9:7. \\[2pt]
\heb{קוֹל דְּמָמָה דַקָּה} \emph{qol demamá daqqá} & \heb{דמם} & ``un sonido de silencio delgado'' (RVR1960: ``un silbo apacible y delicado''). \\
\end{tabular}
\end{filologia}

Dios no está en el viento, ni en el terremoto, ni en el fuego, sino en el silbo apacible (19:11--12). Elías repite su queja palabra por palabra (19:14), y Dios lo corrige: ``yo haré que queden en Israel siete mil, cuyas rodillas no se doblaron ante Baal'' (19:18). El celo de Elías era verdadero, pero había terminado en soledad y desaliento. El celo de Dios, en cambio, guarda un remanente. El contraste enseña que el celo humano por Dios necesita ser corregido por el celo de Dios mismo.

%===================================================================
\section{El celo en el judaísmo}

\begin{nota}[Precisión]
Esta sección expone interpretaciones judías antiguas y medievales. Los textos rabínicos se citan desde Sefaria (Mekhilta, edición Beeri; Talmud, edición William Davidson; Rashi, edición vocalizada; Targum; Maimónides, traducción hebrea de Ibn Tibbon). Donde no se indica otra cosa, la traducción es propia.
\end{nota}

\subsection{``Yo domino el celo''}
La \emph{Mekhilta de Rabí Ismael}, el comentario rabínico antiguo al Éxodo, se detiene en ``Dios celoso'' (Éx 20:5) y hace una pregunta audaz: ``¿Acaso soy un Dios de celo?''. Responde:

\hebverso{אֲנִי שַׁלִּיט בַּקִּנְאָה וְאֵין הַקִּנְאָה שׁוֹלֶטֶת בִּי}

``Yo domino el celo, y el celo no me domina a mí; yo domino el sueño, y el sueño no me domina a mí, como está dicho: `He aquí, no se adormecerá ni dormirá el que guarda a Israel' (Sal 121:4)'' (\emph{Mekhilta}, Bajódesh 6). Y añade otra explicación:

\hebverso{בְּקִנְאָה אֲנִי נִפְרָע מֵהֶם מִן עֲבוֹדָה זָרָה \\ אֲבָל רַחוּם וְחַנּוּן אֲנִי בִּדְבָרִים אֲחֵרִים}

``Con celo castigo la idolatría; pero en las demás cosas soy misericordioso y clemente'' (ibíd.). Los sabios llegan por su camino a la conclusión del material de estudio: el celo de Dios no es una pasión descontrolada; es soberano y apunta específicamente a la idolatría.

\subsection{Rabán Gamaliel y los filósofos}
El Talmud recoge diálogos de Rabán Gamaliel con interlocutores paganos sobre Dt 4:24, ``fuego consumidor, Dios celoso''. Un filósofo pregunta por qué Dios se enoja con los idólatras y no destruye al ídolo. Gamaliel responde con una parábola: un hijo pone a su perro el nombre de su padre y jura ``por la vida del perro, mi padre''. ¿Con quién se enoja el rey, con el hijo o con el perro? Con el hijo (\emph{Avodá Zará} 54b). El ídolo no es rival: es un perro con nombre robado. En otra parábola, un rey que sale a la guerra pelea con los vivos, no con los muertos (ibíd.; \emph{Mekhilta}, Bajódesh 6).

El general Agripas plantea la objeción más aguda: ``¿Acaso no se tiene celo solo de un igual? El sabio del sabio, el fuerte del fuerte, el rico del rico''. Gamaliel responde:

\hebverso{לְאָדָם שֶׁנָּשָׂא אִשָּׁה עַל אִשְׁתּוֹ}

``Es como un hombre que toma una segunda esposa además de la suya: si es más distinguida que ella, la primera no tiene celos; si es inferior, sí los tiene'' (\emph{Avodá Zará} 55a). La parábola invierte la lógica de los celos humanos. El celo de Dios no prueba que el ídolo esté a su altura; prueba que es indigno, ``lo que no es Dios'' (Dt 32:21).

\subsection{Rashi, Ónkelos y Maimónides}
Rashi explica ``Dios celoso'' (Éx 20:5) así: \heb{מְקַנֵּא לִפָּרַע וְאֵינוֹ עוֹבֵר עַל מִדָּתוֹ לִמְחֹל עַל עֲבוֹדַת אֱלִילִים}, ``es celoso para castigar y no pasa por alto su medida para perdonar la idolatría''; y añade que todo \emph{qanná} significa ``dar el corazón a castigar''. Sobre ``cuyo nombre es Celoso'' (Éx 34:14) dice: \heb{אוֹחֵז בְּנִצְחוֹנוֹ וּפוֹרֵעַ מֵאוֹיְבָיו}, ``se aferra a su victoria y castiga a sus enemigos''.

El Targum Ónkelos conserva \emph{El qanná} en los tres textos, pero en Dt 4:24 introduce un matiz: \heb{מֵימְרֵיהּ אֶשָּׁא אָכְלָא הוּא אֵל קַנָּא}, ``\emph{su Palabra} (\emph{Memrá}) es fuego consumidor, Dios celoso''. El Targum evita decir que Dios mismo es fuego.

Maimónides observa una regularidad en toda la Biblia: ``Debes saber que, si examinas toda la Torá y todos los libros de los profetas, no encontrarás la expresión `ira', ni `enojo', ni `\emph{celo}' sino en relación con la idolatría'' (\emph{Guía de los perplejos} I:36). Por eso, añade, solo al idólatra se lo llama ``enemigo'' o ``aborrecedor'' de Dios.

\begin{nota}[Precisión]
Maimónides agrega que creer que Dios es corpóreo provoca su celo más gravemente que la idolatría misma (I:36). Es una lectura filosófica medieval.
\end{nota}

\subsection{Pinjás y Elías en la tradición}
La tradición judía reflexionó mucho sobre los dos celosos del relato bíblico. La Misná recoge un caso excepcional: \heb{קַנָּאִין פּוֹגְעִין בּוֹ}, ``los celosos lo hieren'' (\emph{Sanedrín} 9:6), que el Talmud rodea de condiciones y remite al ejemplo de Pinjás (\emph{Sanedrín} 82a). El \emph{Sifré} a Números llama a Pinjás ``celoso hijo de celoso'' (\heb{קַנָּאִי בֶּן קַנָּאִי}); sobre ``tuvo celo por su Dios'' cita Is 53:12, ``por cuanto derramó su vida hasta la muerte'', y lee ``e hizo expiación'' como algo que sigue haciendo ``hasta que resuciten los muertos'' (\emph{Sifré Bemidbar} 131).

El Targum Pseudo-Jonatán identifica a los dos personajes: \heb{פִּינְחָס הוּא אֵלִיָּהוּ}, ``Pinjás es Elías'', el sumo sacerdote que será enviado a los exiliados al final de los días (a Éx 6:18); y en Nm 25:12 dice que Dios lo hará ``ángel del pacto'' que anunciará la redención (cf. Mal 3:1). Los \emph{Pirqé de Rabí Eliézer} recogen la misma identificación (cap. 47) y ponen en boca de Dios una corrección a Elías en Horeb: ``Siempre estás celando'' (cap. 29).

En el período del Segundo Templo, el celo de Pinjás se volvió modelo. Matatías, al iniciar la revuelta macabea, ``ardió de celo por la ley, como lo hizo Pinjás'', y en su testamento recuerda que ``Pinjás, nuestro padre, por su ardiente celo recibió el pacto de un sacerdocio eterno'' (1 Mac 2:26, 54). En Qumrán, la \emph{Regla de la comunidad} pide que el maestro sea ``hombre celoso por el precepto'', a la espera del día de la venganza (1QS IX, 23; García Martínez, 1992).

\begin{nota}[Precisión]
1 Macabeos se cita según la LXX de Rahlfs-Hanhart, con traducción propia. La línea de 1QS se cita según la edición de García Martínez (1992); conviene cotejarla en la Leon Levy Dead Sea Scrolls Digital Library.
\end{nota}

%===================================================================
\section{Isaías 9: el celo que envía al Mesías}

\subsection{Contexto: Acaz y la crisis siria}
Hacia el año 734 a.\,e.\,c., Siria y el reino del norte (Efraín) atacan Jerusalén para forzar a Judá a unirse contra Asiria (Is 7:1--2). El rey de la casa de David, Acaz, rehúsa la señal que Dios le ofrece (7:12) y busca la protección de Asiria (2 R 16:7). Isaías anuncia la señal de Emanuel (7:14), describe un país en ``tribulación y tinieblas'' (8:22) y, de pronto, anuncia la luz:

\begin{cita}El pueblo que andaba en tinieblas vio gran luz; los que moraban en tierra de sombra de muerte, luz resplandeció sobre ellos.\fuente{(Is 9:2 [9:1])}\end{cita}

Humanamente, la promesa hecha a David (``tu trono será estable eternamente'', 2 S 7:16) está en peligro: el heredero es infiel y el reino, pequeño. Justo ahí Isaías anuncia un hijo de David con un reino sin límite.

\subsection{El texto}
\hebverso{כִּי־יֶלֶד יֻלַּד־לָנוּ בֵּן נִתַּן־לָנוּ וַתְּהִי הַמִּשְׂרָה עַל־שִׁכְמוֹ \\ וַיִּקְרָא שְׁמוֹ פֶּלֶא יוֹעֵץ אֵל גִּבּוֹר אֲבִיעַד שַׂר־שָׁלוֹם \\ לְםַרְבֵּה הַמִּשְׂרָה וּלְשָׁלוֹם אֵין־קֵץ עַל־כִּסֵּא דָוִד וְעַל־מַמְלַכְתּוֹ \\ קִנְאַת יְהוָה צְבָאוֹת תַּעֲשֶׂה־זֹּאת}
\begin{cita}Porque un niño nos es nacido, hijo nos es dado, y el principado sobre su hombro; y se llamará su nombre Admirable, Consejero, Dios Fuerte, Padre Eterno, Príncipe de Paz. Lo dilatado de su imperio y la paz no tendrán límite, sobre el trono de David y sobre su reino, disponiéndolo y confirmándolo en juicio y en justicia desde ahora y para siempre. El celo de \SENOR{} de los ejércitos hará esto.\fuente{(Is 9:6--7 [9:5--6])}\end{cita}

\begin{filologia}[Is 9:6--7 (9:5--6)]
\begin{tabular}{@{}L{3.2cm}L{1.6cm}L{10cm}@{}}
\heb{יֶלֶד יֻלַּד} / \heb{בֵּן נִתַּן} & \heb{ילד}, \heb{נתן} & ``un niño nos ha nacido, un hijo nos ha sido dado'': nacido (humano) y dado (don de Dios). \\[2pt]
\heb{הַמִּשְׂרָה} \emph{hammisrá} & \heb{שׂרר} & ``el principado'', el gobierno: sobre su hombro, como la llave de David (Is 22:22). \\[2pt]
\heb{אֵל גִּבּוֹר} \emph{El gibbor} & \heb{גבר} & ``Dios Fuerte'': en Is 10:21 el mismo título designa a Dios (``el remanente volverá al Dios fuerte''). \\[2pt]
\heb{לְםַרְבֵּה} \emph{lemarbé} & \heb{רבה} & ``para multiplicar''. El texto masorético escribe una \emph{mem} final (\heb{ם}) en medio de la palabra. \\[2pt]
\heb{קִנְאַת יְהוָה צְבָאוֹת} & \heb{קנא} & ``el celo del \SENOR{} de los ejércitos'': \emph{qin'at}, sustantivo en estado constructo. \\[2pt]
\heb{תַּעֲשֶׂה־זֹּאת} \emph{ta'asé zot} & \heb{עשׂה} & ``hará esto'': verbo femenino, que concuerda con \emph{qin'á}. El celo mismo es el sujeto que actúa. \\
\end{tabular}
\end{filologia}

\subsection{El celo como sujeto}
La frase final cierra el oráculo y lo garantiza. Isaías no dice ``el poder'' ni ``la sabiduría'' del \SENOR, sino su celo. La gramática lo subraya: el verbo femenino \heb{תַּעֲשֶׂה} concuerda con \emph{qin'á}; el celo es el agente. La venida del rey prometido nace del amor ardiente de Dios por su pueblo y por su promesa. El título ``de los ejércitos'' es el mismo con que Elías dijo haber celado por Dios (1 R 19:10); en Isaías es Dios quien cela.

\begin{center}\small\renewcommand{\arraystretch}{1.2}
\begin{tabular}{@{}L{2.8cm}L{5.4cm}L{5.4cm}@{}}
\toprule
\textbf{Testigo} & \textbf{Texto} & \textbf{Traducción} \\ \midrule
TM (9:6) & \heb{קִנְאַת יְהוָה צְבָאוֹת תַּעֲשֶׂה־זֹּאת} & ``el celo del \SENOR{} de los ejércitos hará esto'' \\
LXX (9:6) & \grc{ὁ ζῆλος κυρίου σαβαωθ ποιήσει ταῦτα} & ``el celo del Señor Sabaot hará estas cosas'' \\
Targum Jonatán (9:6) & \heb{בְּמֵימְרָא דַייָ צְבָאוֹת תִּתְעֲבֵיד דָא} & ``por la Palabra del \SENOR{} de los ejércitos se hará esto'' \\
\bottomrule
\end{tabular}\par
{\footnotesize\color{gris}\textit{Tabla 9.} Is 9:7 [9:6] en los testigos antiguos. \textit{Nota.} BHS; Rahlfs-Hanhart (2006); Targum Jonatán (Sefaria).}
\end{center}

\begin{nota}[Precisión textual]
El Targum sustituye ``celo'' por \emph{Memrá}, ``Palabra'', como hace en otros lugares para no atribuir pasiones a Dios; la Septuaginta conserva \grc{ζῆλος}. El texto hebreo, que sigue la traducción, dice ``celo''.
\end{nota}

\subsection{Tres veces la misma frase}
La fórmula \heb{קִנְאַת יְהוָה צְבָאוֹת תַּעֲשֶׂה־זֹּאת} aparece tres veces en la Biblia: aquí, en Is 37:32 y en su paralelo de 2 R 19:31. Las otras dos pertenecen al oráculo sobre el asedio de Senaquerib en tiempos de Ezequías:

\begin{cita}Porque de Jerusalén saldrá un remanente, y del monte de Sion los que se salven. El celo de \SENOR{} de los ejércitos hará esto.\fuente{(Is 37:32)}\end{cita}
\begin{cita}Porque yo ampararé a esta ciudad para salvarla, por amor de mí mismo, y por amor de David mi siervo.\fuente{(Is 37:35)}\end{cita}

El celo salva a Jerusalén ``por amor de mí mismo, y por amor de David''. El mismo celo que preserva un remanente en la crisis asiria es el que traerá al Hijo de David.

\begin{nota}[Precisión textual]
En 2 R 19:31 el códice escribe solo ``el celo del \SENOR'', y los masoretas indican leer ``de los ejércitos'' (\emph{qeré} sin \emph{ketiv}). La conexión teológica entre los tres textos es lectura propia de esta serie.
\end{nota}

\subsection{La \emph{mem} cerrada}
La \emph{mem} final escrita en medio de \heb{לְםַרְבֵּה} llamó la atención de los sabios. El Talmud la explica así:

\hebverso{בִּיקֵּשׁ הַקָּדוֹשׁ בָּרוּךְ הוּא לַעֲשׂוֹת חִזְקִיָּהוּ מָשִׁיחַ וְסַנְחֵרִיב גּוֹג וּמָגוֹג}

``¿Por qué toda \emph{mem} en medio de palabra es abierta, y esta es cerrada? El Santo, bendito sea, quiso hacer a Ezequías Mesías y a Senaquerib Gog y Magog''. La medida de la justicia objetó que Ezequías no había cantado alabanza después de tantos milagros; ``por eso se cerró'' (\emph{Sanedrín} 94a). Los sabios unen así, por su camino, Is 9:7 con la liberación de Jerusalén frente a Senaquerib, el mismo escenario de Is 37:32.

Rashi, que lee el niño como Ezequías, comenta la frase final: ``el celo del \SENOR{} de los ejércitos'', que tuvo celo por Sion frente a lo que tramaron Aram y Pecaj; y ``hará esto'',

\hebverso{וְלֹא אָחָז רָאוּי לְכָךְ וְאַף זְכוּת אָבוֹת תַּמָּה}

``y Acaz no es digno de ello, y aun el mérito de los padres se ha agotado'' (Rashi a Is 9:6). La observación es valiosa para la teología del atributo: el celo de Dios actúa donde no hay mérito. Es pura gracia.

\begin{nota}[Precisión]
La \emph{mem} cerrada es un dato del texto masorético; su sentido es interpretación rabínica. Rashi sigue la lectura histórica medieval (Ezequías) y discute en otros lugares de su comentario la lectura cristiana.
\end{nota}

\subsection{¿Quién es el niño?}
La tradición judía antigua leyó el texto como mesiánico. El Targum Jonatán traduce el final de Is 9:5 [heb.]: \heb{מְשִׁיחָא דִשְׁלָמָא יַסְגֵי עֲלָנָא בְּיוֹמוֹהִי}, ``el Mesías, en cuyos días abundará la paz sobre nosotros''. Rashi, en cambio, lo aplica a Ezequías. El Nuevo Testamento lo lee cumplido en Jesús: el ángel anuncia que ``el Señor Dios le dará el trono de David su padre; y reinará sobre la casa de Jacob para siempre, y su reino no tendrá fin'' (Lc 1:32--33), y Mateo cita Is 9:1--2 al comienzo del ministerio de Jesús en Galilea (Mt 4:15--16).

Los nombres del niño superan a cualquier rey de Judá. ``Dios Fuerte'' es, en el capítulo siguiente, Dios mismo (Is 10:21); ``Padre Eterno'' y un reino ``desde ahora y para siempre'' exceden la vida de Ezequías. El Targum ya lo leía como el Mesías; el Nuevo Testamento identifica a ese Mesías.

\subsection{Por qué el celo envía al Mesías}
El encadenamiento de Isaías es claro: hay una promesa (2 S 7:16), un heredero infiel (Acaz), un pueblo en tinieblas y una amenaza extranjera. Lo que hace avanzar la historia no es el mérito de Israel, sino el celo de Dios. El celo es el amor que ``no suelta'' lo que es suyo (Cnt 8:6). Por eso no solo destruye ídolos: también rescata. Frente a un rey que busca a Asiria y a un pueblo que anda en tinieblas, Dios no permitirá que otro rey, otro dios u otra potencia se quede con su pueblo. Envía al Rey.

En Isaías 9 el celo no es fuego contra Israel, sino ardor a favor de Israel. El Mesías no viene porque Israel lo mereciera, sino porque Dios es celoso.

\begin{nota}[Precisión]
Síntesis teológica de esta serie, a partir de los textos de Isaías y de la observación de Rashi.
\end{nota}

%===================================================================
\section{El celo de Dios en Cristo, leído con Pablo}

\subsection{``El celo de tu casa me consume''}
Al purificar el templo, Jesús encarna el celo de Dios por su casa:

\begin{cita}Y haciendo un azote de cuerdas, echó fuera del templo a todos, y las ovejas y los bueyes; y esparció las monedas de los cambistas, y volcó las mesas; y dijo a los que vendían palomas: Quitad de aquí esto, y no hagáis de la casa de mi Padre casa de mercado. Entonces se acordaron sus discípulos que está escrito: El celo de tu casa me consume.\fuente{(Jn 2:15--17)}\end{cita}

\hebverso{כִּי־קִנְאַת בֵּיתְךָ אֲכָלָתְנִי וְחֶרְפּוֹת חוֹרְפֶיךָ נָפְלוּ עָלָי}
\begin{cita}Porque me consumió el celo de tu casa; Y los denuestos de los que te vituperaban cayeron sobre mí.\fuente{(Sal 69:9 [69:10])}\end{cita}

Juan cita la primera mitad del verso. Pablo cita la segunda y la aplica a Cristo: ``Porque ni aun Cristo se agradó a sí mismo; antes bien, como está escrito: Los vituperios de los que te vituperaban, cayeron sobre mí'' (Ro 15:3). Leídas juntas, las dos citas muestran a Cristo ardiendo de celo por la casa de Dios y cargando el oprobio de los que ofenden a Dios.

\begin{nota}[Precisión]
Juan cambia el pasado del Salmo en la LXX (\grc{κατέφαγεν}) por un futuro (\grc{καταφάγεταί}), ``me consumirá'': el celo llevará a Jesús a la muerte (Carson, 1991).
\end{nota}

\subsection{``Os celo con celo de Dios''}
Pablo hereda el lenguaje matrimonial de los profetas y lo aplica a la iglesia:

\begin{cita}Porque os celo con celo de Dios; pues os he desposado con un solo esposo, para presentaros como una virgen pura a Cristo. Pero temo que como la serpiente con su astucia engañó a Eva, vuestros sentidos sean de alguna manera extraviados de la sincera fidelidad a Cristo.\fuente{(2 Co 11:2--3)}\end{cita}

\begin{filologia}[2 Co 11:2]
\begin{tabular}{@{}L{3.2cm}L{1.6cm}L{10cm}@{}}
\grc{ζηλῶ γὰρ ὑμᾶς θεοῦ ζήλῳ} & \grc{ζῆλος} & ``os celo con celo \emph{de Dios}'': como Pinjás, Pablo comparte el celo de Dios. \\[2pt]
\grc{ἡρμοσάμην} & \grc{ἁρμόζω} & ``os comprometí en matrimonio'': Pablo actúa como el amigo del novio. \\[2pt]
\grc{ἑνὶ ἀνδρί} & --- & ``a un solo esposo'': la exclusividad del pacto. \\
\end{tabular}
\end{filologia}

En los profetas, el Esposo de Israel es el \SENOR: ``tu marido es tu Hacedor; \SENOR{} de los ejércitos es su nombre'' (Is 54:5; cf. Os 2:19--20). En Pablo, el Esposo es Cristo. El celo que el Sinaí atribuía a Dios por Israel se ejerce ahora por la iglesia en favor de Cristo.

\subsection{Provocar a celos: Deuteronomio 32 en Romanos}
Pablo encuentra en el cántico de Moisés la clave del destino de Israel:

\begin{center}\small\renewcommand{\arraystretch}{1.18}
\begin{tabular}{@{}L{2.0cm}L{6.0cm}L{5.8cm}@{}}
\toprule
\textbf{Texto} & \textbf{Original} & \textbf{RVR1960} \\ \midrule
Dt 32:21 (TM) & \heb{הֵם קִנְאוּנִי בְלֹא־אֵל} / \heb{וַאֲנִי אַקְנִיאֵם בְּלֹא־עָם} & ``me movieron a celos con lo que no es Dios\ldots{} yo también los moveré a celos con un pueblo que no es pueblo'' \\
Dt 32:21 (LXX) & \grc{αὐτοὶ παρεζήλωσάν με ἐπ’ οὐ θεῷ\ldots{} κἀγὼ παραζηλώσω αὐτοὺς ἐπ’ οὐκ ἔθνει} & --- \\
Ro 10:19 & \grc{ἐγὼ παραζηλώσω ὑμᾶς ἐπ’ οὐκ ἔθνει} & ``Yo os provocaré a celos con un pueblo que no es pueblo'' \\
\bottomrule
\end{tabular}\par
{\footnotesize\color{gris}\textit{Tabla 10.} Dt 32:21 en Romanos. \textit{Nota.} BHS; Rahlfs-Hanhart (2006); NA28.}
\end{center}

\begin{cita}Digo, pues: ¿Han tropezado los de Israel para que cayesen? En ninguna manera; pero por su transgresión vino la salvación a los gentiles, para provocarles a celos.\fuente{(Ro 11:11)}\end{cita}
\begin{cita}por si en alguna manera pueda provocar a celos a los de mi sangre, y hacer salvos a algunos de ellos.\fuente{(Ro 11:14)}\end{cita}

La lógica es la del cántico: Israel provocó a Dios con un ``no-dios''; Dios provoca a Israel con un ``no-pueblo''. Pero Pablo descubre en ello un plan de misericordia: la salvación llega a las naciones para despertar a Israel, ``y luego todo Israel será salvo'' (Ro 11:26). El celo de Dios no termina en descarte, sino en salvación, y Pablo concluye con la doxología de Ro 11:33: ``¡Oh profundidad de las riquezas de la sabiduría y de la ciencia de Dios!''.

\subsection{``¿Provocaremos a celos al Señor?''}
En 1 Corintios 10, Pablo advierte a los creyentes que participan en banquetes idolátricos:

\begin{cita}Antes digo que lo que los gentiles sacrifican, a los demonios lo sacrifican, y no a Dios; y no quiero que vosotros os hagáis partícipes con los demonios. No podéis beber la copa del Señor, y la copa de los demonios; no podéis participar de la mesa del Señor, y de la mesa de los demonios. ¿O provocaremos a celos al Señor? ¿Somos más fuertes que él?\fuente{(1 Co 10:20--22)}\end{cita}

``A los demonios\ldots{} y no a Dios'' cita Dt 32:17 en la LXX (\grc{ἔθυσαν δαιμονίοις καὶ οὐ θεῷ}), y ``¿provocaremos a celos?'' (\grc{παραζηλοῦμεν}) es el verbo de Dt 32:21. El ``Señor'' es el de ``la copa del Señor'' y ``la mesa del Señor'': Jesucristo. Pablo pone a Cristo en el lugar del Dios celoso del cántico de Moisés.

\begin{nota}[Precisión]
La lectura cristológica de ``Señor'' en 1 Co 10:22 la sostienen Fee (1987) y Hays (1989); es la más probable por el contexto de la Cena del Señor (10:16--21).
\end{nota}

\subsection{El celo de Saulo}
Pablo conocía el celo desde dentro:

\begin{cita}y en el judaísmo aventajaba a muchos de mis contemporáneos en mi nación, siendo mucho más celoso de las tradiciones de mis padres.\fuente{(Gá 1:14)}\end{cita}
\begin{cita}en cuanto a celo, perseguidor de la iglesia; en cuanto a la justicia que es en la ley, irreprensible. Pero cuantas cosas eran para mí ganancia, las he estimado como pérdida por amor de Cristo.\fuente{(Flp 3:6--7)}\end{cita}

Pablo se llama \grc{ζηλωτής}, ``celoso'' (Gá 1:14; Hch 22:3): se ubica en la línea de Pinjás, Elías y los Macabeos. Ese celo lo llevó a perseguir a la iglesia. De Israel dirá: ``tienen celo de Dios, pero no conforme a ciencia'' (Ro 10:2, \grc{οὐ κατ’ ἐπίγνωσιν}). En Cristo, el celo no desaparece: cambia de objeto. El perseguidor ahora cela por la iglesia ``con celo de Dios'' (2 Co 11:2).

\begin{nota}[Precisión]
El trasfondo de Pinjás en el celo de Pablo es tesis de Hengel (1989) y de Wright (2013); Dunn (1998) la recoge. El paralelo entre Sal 106:31 y Gn 15:6, ya notado, sugiere por qué Pablo, en Ro 4, funda la justicia en la fe de Abraham y no en el celo.
\end{nota}

\subsection{Un pueblo celoso de buenas obras}
El Dios celoso redime un pueblo que le pertenece, y ese pueblo aprende a celar por el bien:

\begin{cita}quien se dio a sí mismo por nosotros para redimirnos de toda iniquidad y purificar para sí un pueblo propio, celoso de buenas obras.\fuente{(Tit 2:14)}\end{cita}

``Pueblo propio'' traduce \grc{λαὸν περιούσιον}, la expresión de la LXX para la \emph{segul.lá} de Éx 19:5; ``celoso de buenas obras'' es \grc{ζηλωτὴν καλῶν ἔργων}. El celo de Dios crea un pueblo celoso.

\begin{nota}[Precisión]
La autoría paulina de Tito se discute. La idea tiene respaldo en cartas indiscutidas: ``Bueno es mostrar celo en lo bueno siempre'' (Gá 4:18).
\end{nota}

\subsection{El celo de Dios, cumplido en Cristo}
Leído con Pablo, el celo del Sinaí se cumple en Cristo en cuatro momentos: el celo de Dios \emph{envía} al Hijo (Is 9:7); el celo por la casa de Dios \emph{consume} al Hijo (Jn 2:17; Ro 15:3); el celo de Dios \emph{guarda} a la iglesia como esposa (2 Co 11:2); el celo de Dios \emph{salva} a Israel y a las naciones (Ro 11:11--26). El Dios que no da su gloria a otro (Is 42:8) la muestra ``en la faz de Jesucristo'' (2 Co 4:6), porque ``en él habita corporalmente toda la plenitud de la Deidad'' (Col 2:9). Quien ve el celo de Cristo por la casa de su Padre ve el celo del \SENOR.

%===================================================================
\section{Aplicación: ¿cómo me bendice el celo de Dios?}

La cuarta pregunta del material es la más personal. Su respuesta: el celo de Dios me bendice porque me rescata de mis ídolos, me protege de la autodestrucción, preserva el pacto, me disciplina para restaurarme y garantiza mi salvación. ``Sin el celo de Dios, pereceríamos''.

\begin{cita}Porque yo \SENOR, Dios tuyo, el Santo de Israel, soy tu Salvador; a Egipto he dado por tu rescate, a Etiopía y a Seba por ti. Porque a mis ojos fuiste de gran estima, fuiste honorable, y yo te amé; daré, pues, hombres por ti, y naciones por tu vida.\fuente{(Is 43:3--4)}\end{cita}

Cuatro respuestas se desprenden del estudio:

\begin{enumerate}
\item \textbf{Examinar el corazón.} El celo de Dios se provoca cuando un ``amante'' ocupa su lugar. La pregunta no es si tengo ídolos de piedra, sino qué compite hoy con Dios: ``El que ama a padre o madre más que a mí, no es digno de mí'' (Mt 10:37).
\item \textbf{Huir de la idolatría.} Pablo no manda negociar con los ídolos, sino huir: ``Por tanto, amados míos, huid de la idolatría'' (1 Co 10:14). Juan cierra su carta igual: ``Hijitos, guardaos de los ídolos'' (1 Jn 5:21). Santiago usa el lenguaje de Oseas: ``¡Oh almas adúlteras! ¿No sabéis que la amistad del mundo es enemistad contra Dios?'' (Stg 4:4).
\item \textbf{Recibir la disciplina como amor.} ``Yo reprendo y castigo a todos los que amo; sé, pues, celoso, y arrepiéntete'' (Ap 3:19). El ídolo promete vida y entrega muerte; el celo de Dios lo destruye antes de que el ídolo nos destruya.
\item \textbf{Descansar en su celo.} ``El celo de \SENOR{} de los ejércitos hará esto'' (Is 9:7). La última palabra no la tiene mi fidelidad, sino su celo.
\end{enumerate}

Como reflejo, el celo es comunicable: el creyente arde por la gloria de Dios y por la pureza de la iglesia (Gá 4:18; Tit 2:14; Ro 12:11). Pero el celo humano necesita ``ciencia'' (Ro 10:2), está siempre expuesto a volverse envidia o ``celos amargos'' (Stg 3:14--16; Gá 5:20) y requiere la corrección que Dios dio a Elías en Horeb. El buen celo se parece al de Pinjás, que celó con el celo \emph{de Dios}, y al de Pablo, que celó ``con celo de Dios'' por una iglesia que no era suya, sino de Cristo.

%===================================================================
\section{Síntesis}

``\SENOR, cuyo nombre es Celoso, Dios celoso es'' (Éx 34:14). En ese verso, la misma palabra, \emph{qanná}, es nombre y adjetivo; el celo pertenece a la identidad revelada de Dios. Es el ardor de su santo amor de pacto: no tolera rivales porque el rival es ``lo que no es Dios'' y destruye a su pueblo. Se enciende con la idolatría, se manifiesta en disciplina y juicio, y arde \emph{por} Sion, por su Nombre y por su promesa. Por ese celo vino el Hijo de David (Is 9:7), ardió por la casa de su Padre (Jn 2:17), guarda a la iglesia como esposa (2 Co 11:2) y salvará a Israel y a las naciones (Ro 11:11--26).

\begin{center}\small\renewcommand{\arraystretch}{1.18}
\begin{tabular}{@{}L{2.4cm}L{11.4cm}@{}}
\toprule
\textbf{Mirada} & \textbf{Aporte al atributo} \\ \midrule
Hebrea & \emph{qanná} como nombre y adjetivo (Éx 34:14), solo de Dios; ardor; celo \emph{por} (\heb{לְ}) frente a envidia (\heb{בְּ}); el pacto como matrimonio; Pinjás y Elías; Is 9:7. \\
Judía & ``Yo domino el celo'' (Mekhilta); solo contra la idolatría (Maimónides); la segunda esposa (\emph{Avodá Zará} 55a); la \emph{mem} cerrada (\emph{Sanedrín} 94a); ``Acaz no es digno'' (Rashi). \\
Pink & ``Nuestro Dios es fuego consumidor'': reverencia; el amor de Dios no se modela según la emoción humana. \\
Paulina & ``Celo de Dios'' por la iglesia; Dt 32:21 en Ro 10--11 y 1 Co 10:22; el celo de Saulo transformado; un pueblo celoso de buenas obras. \\
\bottomrule
\end{tabular}\end{center}

\begin{cita}Mas alábese en esto el que se hubiere de alabar: en entenderme y conocerme, que yo soy \SENOR, que hago misericordia, juicio y justicia en la tierra; porque estas cosas quiero, dice \SENOR.\fuente{(Jer 9:24 [9:23])}\end{cita}

Conocer al Dios celoso es saber que no nos comparte con nadie. Lo que en los celos humanos sería amenaza, en Dios es seguridad: el que nos reclama para sí es el que no permitirá que otro se quede con nosotros.

%===================================================================
\section*{Glosario}
\begin{description}\setlength{\itemsep}{2pt}
\item[\emph{El gibbor}] (\heb{אֵל גִּבּוֹר}). ``Dios Fuerte'' (Is 9:6; 10:21).
\item[\emph{El qanná}] (\heb{אֵל קַנָּא}). ``Dios celoso'' (Éx 20:5; 34:14; Dt 4:24).
\item[\emph{Hiqní}] (\heb{הִקְנִיא}). Provocar a celos (Dt 32:16, 21).
\item[\emph{Jemá}] (\heb{חֵמָה}). Calor, furor; acompaña a \emph{qin'á}.
\item[\emph{Kanná'im}] (\heb{קַנָּאִים}). ``Los celosos'', zelotes (Misná \emph{Sanedrín} 9:6).
\item[\emph{Memrá}] (\heb{מֵימְרָא}). ``Palabra'' del \SENOR{} en los targumes.
\item[\emph{Parazēlóō}] (\grc{παραζηλόω}). Provocar a celos (Ro 10:19; 11:11, 14; 1 Co 10:22).
\item[\emph{Qanná}] (\heb{קַנָּא}). Celoso; adjetivo reservado a Dios.
\item[\emph{Qannó}] (\heb{קַנּוֹא}). Celoso (Jos 24:19; Nah 1:2).
\item[\emph{Qinné}] (\heb{קִנֵּא}). Tener celo por (con \heb{לְ}); envidiar (con \heb{בְּ}).
\item[\emph{Qin'á}] (\heb{קִנְאָה}). Celo, ardor; celos; envidia.
\item[\emph{Segul.lá}] (\heb{סְגֻלָּה}). Tesoro personal (Éx 19:5).
\item[\emph{Sémel haqqin'á}] (\heb{סֵמֶל הַקִּנְאָה}). ``La imagen del celo'' (Ez 8:3, 5).
\item[\emph{Shalhévetyá}] (\heb{שַׁלְהֶבֶתְיָה}). ``Llama de Yah'' (Cnt 8:6).
\item[\emph{Zaná}] (\heb{זָנָה}). Fornicar; ser infiel al pacto (Éx 34:15).
\item[\emph{Zēlos}] (\grc{ζῆλος}). Celo, ardor; de \grc{ζέω}, ``hervir''.
\item[\emph{Zēlōtēs}] (\grc{ζηλωτής}). Celoso; zelote (Gá 1:14; Tit 2:14).
\end{description}

\section*{Referencias}
\refapa{Bauer, W., Danker, F.~W., Arndt, W.~F., \& Gingrich, F.~W. (2000). \emph{A Greek-English lexicon of the New Testament and other early Christian literature} (3.ª ed.). University of Chicago Press.}
\refapa{Brown, F., Driver, S.~R., \& Briggs, C.~A. (1906). \emph{A Hebrew and English lexicon of the Old Testament}. Clarendon Press.}
\refapa{Carson, D.~A. (1991). \emph{The Gospel according to John}. Eerdmans.}
\refapa{Carson, D.~A. (1996). \emph{Exegetical fallacies} (2.ª ed.). Baker.}
\refapa{Chávez, M. (1992). \emph{Diccionario de hebreo bíblico}. Editorial Mundo Hispano.}
\refapa{Dunn, J.~D.~G. (1998). \emph{The theology of Paul the apostle}. Eerdmans.}
\refapa{Elliger, K., \& Rudolph, W. (Eds.). (1997). \emph{Biblia Hebraica Stuttgartensia} (5.ª ed.). Deutsche Bibelgesellschaft.}
\refapa{Fee, G.~D. (1987). \emph{The first epistle to the Corinthians}. Eerdmans.}
\refapa{García Martínez, F. (1992). \emph{Textos de Qumrán}. Trotta.}
\refapa{Hays, R.~B. (1989). \emph{Echoes of Scripture in the letters of Paul}. Yale University Press.}
\refapa{Hengel, M. (1989). \emph{The Zealots} (D. Smith, Trad.). T\&T Clark.}
\refapa{Koehler, L., \& Baumgartner, W. (2001). \emph{The Hebrew and Aramaic lexicon of the Old Testament} (M.~E.~J. Richardson, Trad.; Ed. de estudio). Brill.}
\refapa{Maimónides. (s.f.). \emph{Guía de los perplejos} [I:36] (S. ibn Tibbon, Trad. hebrea). Sefaria. https://www.sefaria.org}
\refapa{\emph{Mekhilta de Rabí Ismael} [Bajódesh 6]. (2019). (Ed. Beeri). Sefaria. https://www.sefaria.org}
\refapa{Nestle, E., \& Aland, K. (Eds.). (2012). \emph{Novum Testamentum Graece} (28.ª ed.). Deutsche Bibelgesellschaft.}
\refapa{Pink, A.~W. (2020). \emph{Los atributos de Dios} (E. Rodríguez Cruz, Trad.). Chapel Library. (Obra original publicada en 1930)}
\refapa{\emph{Pirqé de Rabí Eliézer} [29; 47]. (s.f.). Sefaria. https://www.sefaria.org}
\refapa{Rahlfs, A., \& Hanhart, R. (Eds.). (2006). \emph{Septuaginta} (Ed. rev.). Deutsche Bibelgesellschaft.}
\refapa{Rashi. (s.f.). \emph{Comentario a Éxodo e Isaías} [Éx 20:5; 34:14; Is 9:5--6]. Sefaria. https://www.sefaria.org}
\refapa{Real Academia Española. (2014). \emph{Diccionario de la lengua española} (23.ª ed.). [``celo'']}
\refapa{\emph{Santa Biblia: Reina-Valera 1960}. (1960). Sociedades Bíblicas en América Latina.}
\refapa{\emph{Sifré Bemidbar} [131]. (s.f.). Sefaria. https://www.sefaria.org}
\refapa{Steinsaltz, A. (Ed.). (2017). \emph{Koren Talmud Bavli: Edición William Davidson} [Sanedrín 82a, 94a; Avodá Zará 54b--55a]. Sefaria. https://www.sefaria.org}
\refapa{Strong, J. (1890). \emph{The exhaustive concordance of the Bible}. Hunt \& Eaton.}
\refapa{\emph{Targum Jonatán a los Profetas} [Is 9:5--6]; \emph{Targum Ónkelos} [Éx 20:5; 34:14; Dt 4:24]; \emph{Targum Pseudo-Jonatán} [Éx 6:18; Nm 25:12]. (s.f.). Sefaria. https://www.sefaria.org}
\refapa{Vine, W.~E. (1999). \emph{Diccionario expositivo de palabras del Antiguo y del Nuevo Testamento exhaustivo}. Caribe.}
\refapa{Wright, N.~T. (2013). \emph{Paul and the faithfulness of God}. Fortress Press.}

\end{document}
